\chapter{Running a Bond}
\label{ch:running}

\section{Before you start}

\begin{enumerate}[leftmargin=*]
\item The machine has homed successfully and shows no error.
\item A configuration is loaded --- you are on the Parameter Edit page, not the
  configuration list.
\item The workpiece is secured and within reach at the search height.
\item A tail of wire protrudes from the tool.
\end{enumerate}

\begin{wbnote}
The machine only accepts a bond trigger when it is idle. If it is running
something else, or an error is latched, the mouse does nothing at all --- there
is no message. Check the display.
\end{wbnote}

\section{The button rhythm}

In the three automatic modes the pattern is the same for both bonds:

\begin{center}
\textbf{press and hold} $\rightarrow$ machine moves into position $\rightarrow$
\textbf{release} $\rightarrow$ machine descends and welds
\end{center}

The release is the commitment. While you hold the button the head has moved to
its search height but will not descend, so you can take as long as you like to
line up. Let go when you are happy.

Between the two bonds the machine waits at \menuitem{Loop Height}
indefinitely. There is no timeout on your reaction --- take the time you need
to position the second pad.

\section{What happens automatically}

Before any cycle begins, the machine always returns itself to a known state:
force off, clamp closed, head at \menuitem{Reset Height}, and every pending
button press discarded.

\begin{wbnote}
That last point is why a button you pressed while the machine was busy does not
cause a surprise cycle afterwards. Every cycle starts from a fresh, deliberate
press.
\end{wbnote}

\section{Stopping}

There is no pause and no cancel. Once a phase is running it runs to completion
or it faults.

\begin{itemize}[leftmargin=*]
\item \textbf{Between the two bonds}, the machine is simply waiting. It will
  wait forever. Nothing is at risk.
\item \textbf{Mid-weld}, the weld finishes in at most \menuitem{Bond Timeout}.
\item \textbf{If you must stop immediately}, press \keycap{RESET}. The machine
  restarts and re-homes. The wire will be left in whatever state it was in ---
  clear it before bonding again.
\end{itemize}

\begin{wbwarning}
\keycap{RESET} is a restart, not an emergency stop. See Chapter~\ref{ch:safety}.
\end{wbwarning}

\section{Reading the ultrasonic report}

After each weld the machine briefly reports what it measured --- resonant
frequency, Q, and delivered power. Glance at it. A frequency that has drifted
towards the edge of your scan window, or a Q that has fallen since yesterday,
predicts bond problems before they show up in pull tests.

\section{If a cycle fails}

The machine shows which operation failed and why. The four faults and what to
do about them are in Chapter~\ref{ch:messages}.

A failed bond cycle is a warning, not an error: clear the message, fix the
cause, and bond again. Only homing and initialisation lock the machine.

\section{Routine checks}

\TODO{Write the shift routine for this shop: what to check at start of shift
(run \keycap{TEST} and record frequency and Q, verify force calibration,
inspect the tool), how often to re-run \keycap{SETUP}, the pull-test schedule
and its acceptance criteria, and when the tool must be changed.}

\section{Loading wire}

\TODO{Write the wire loading procedure with photographs: threading the spool,
routing through the guides and clamp, setting tension, forming the initial tail
and its correct length, and what to do after a wire break mid-cycle.}
